% Introduction and related work draft for an IEEE Transactions on Smart Grid manuscript.
% Intended use in an IEEEtran main file:
% % Introduction and related work draft for an IEEE Transactions on Smart Grid manuscript.
% Intended use in an IEEEtran main file:
% % Introduction and related work draft for an IEEE Transactions on Smart Grid manuscript.
% Intended use in an IEEEtran main file:
% % Introduction and related work draft for an IEEE Transactions on Smart Grid manuscript.
% Intended use in an IEEEtran main file:
% \input{sections/01_introduction_related_work}

\section{Introduction}
\IEEEPARstart{D}{istribution} system restoration is a time-critical control problem in which operators and automation systems must reconnect loads, reconfigure feeders, and coordinate available switching or controllable load resources after faults and outages. Modern active distribution networks increase the value of automated restoration because they include remotely controllable switches, distributed energy resources, microgrid-capable regions, and increasingly detailed sensing. At the same time, these capabilities expand the cyber-physical attack surface: restoration decisions depend not only on the true electrical state of the feeder, but also on the cyber observation made available to the restoration controller.

Learning-based restoration controllers are attractive in this setting because they can produce fast decisions after training and can naturally represent sequential restoration objectives. Graph reinforcement learning (GRL), in particular, is well matched to distribution feeders because feeder topology, switch state, voltage, branch-flow, and load connectivity are graph-structured quantities. A GRL policy can therefore use node, edge, and global features to map post-outage observations to switching and load-restoration actions. However, the same observation-rich structure that makes GRL useful also creates a security concern: if an attacker corrupts selected entries of the controller observation, the policy may choose a physically poor restoration action even though the underlying feeder has not changed.

False-data injection attacks (FDIAs) have been studied extensively in power-system monitoring and state estimation, where carefully designed measurement corruptions can bypass conventional bad-data detection and bias estimated system states \cite{liu2009false,kosut2011malicious,xie2011false,bi2014graphical}. Those results show that cyber measurements cannot be treated as automatically trustworthy. Yet restoration control differs from state estimation in several important ways. The consequence of an attack is not only an erroneous estimate; it is a changed restoration action that is then applied to the physical feeder. Moreover, the most security-critical observation may not be a traditional voltage or flow measurement. It may be a binary action mask, a topology/status input, or another controller-specific feature that constrains the learned policy's action selection.

This paper studies observation-level FDIAs against learning-based distribution restoration under a strict cyber-physical evaluation rule. The attacker corrupts only the observation received by the restoration controller. The physical feeder, outage scenario, and OpenDSS evaluation model remain the ground truth. For each deterministic outage-edge scenario, we compare the clean controller action, the attacked controller action, and the mitigated action on the same physical feeder state. Attack impact is measured using positive restored-energy degradation relative to the clean baseline, so cases in which an attack accidentally improves the restoration outcome are not counted as damage.

The study focuses on direct observations available to the current restoration controller: action/status masks, topology/status inputs, voltage measurements, branch-flow measurements, and global restoration scalars. This boundary is important. We do not claim direct attacks on load demand, DER availability, or explicit fault-status observations unless those quantities are exposed in the controller observation model. Instead, the goal is to quantify which existing controller-observation channels can create physical degradation when corrupted.

The paper is built on two frozen case-study evidence packages. The IEEE 34-bus study shows that action/status attacks dominate the observed vulnerability. In particular, false-unavailable dispatch-load attacks centered on the `Load\_s890' action/status observation create multi-MWh restored-energy degradation, and coordinated attacks involving the same observation reach mean physical degradation above 3 MWh. The IEEE 123-bus study shows a different vulnerability structure. In the baseline family study, branch-flow hiding is the strongest aggregate attack family, while exhaustive action-mask criticality reveals individual load-control observations whose impact is much larger than the baseline action-mask family average. These cross-feeder differences demonstrate that restoration security cannot be summarized by one generic FDIA family; it requires registry-driven criticality analysis over the actual observations used by the controller.

We also evaluate detection and mitigation. The detection layer includes residual and status-integrity checks, with strict and relaxed profiles used to quantify robustness and bypass risk. The mitigation layer screens candidate restoration actions after detection. In IEEE 34, trusted-reference screening substantially reduces residual degradation for high-impact action/status attacks, while relaxed status assumptions with weak candidate screening can leave large residual losses. In IEEE 123, the completed defense-stress study finds no high-impact detection bypass cases, but identifies detected low-recovery residual cases; advanced mitigation variants based on detector-selected, trusted-reference, and horizon-screened candidate actions remove the residual degradation in those hard cases.

The main contributions of this paper are:
\begin{itemize}
    \item An auditable observation-level FDIA assessment workflow for learning-based distribution restoration, where cyber observations are corrupted but all clean, attacked, and mitigated actions are evaluated on the true OpenDSS feeder state.
    \item A controller-observation attack taxonomy covering action/status masks, topology/status inputs, voltage measurements, branch-flow measurements, and global restoration scalars, with explicit claim boundaries for observations not exposed to the current controller.
    \item Frozen IEEE 34-bus and IEEE 123-bus GRL case studies with deterministic outage-edge replay, restored-energy degradation metrics, damaging-case counts, detection rates, bypass rates, mitigation recovery, and residual-risk accounting.
    \item Cross-feeder evidence that the dominant vulnerability mode is feeder and controller-context dependent: IEEE 34 is driven by dispatch-load action/status observations, whereas IEEE 123 combines branch-flow family effects with high-impact individual action-mask targets.
    \item A defense evaluation showing that strict status-integrity detection and trusted-reference or horizon-screened candidate mitigation can substantially reduce residual restoration loss, while detection alone is not sufficient when mitigation candidate selection remains weak.
\end{itemize}

The remainder of the paper is organized as follows. Section II reviews related work on distribution restoration, learning-based restoration, FDIA security, and adversarial reinforcement learning. Section III introduces the restoration environment, observation model, threat model, and defense model. Section IV describes the deterministic experimental workflow and evidence packages. Sections V and VI present the IEEE 34-bus and IEEE 123-bus results, respectively. Section VII discusses companion controller evidence and the boundary of the current claims. Section VIII concludes and outlines the planned IEEE 8500-bus HPC-scale extension.

\section{Related Work}

\subsection{Distribution System Restoration}

Distribution system restoration has traditionally been formulated as an optimization, search, or expert-system problem in which the objective is to maximize restored load or minimize unserved energy subject to radiality, voltage, thermal, switching, and operational constraints. Classical feeder restoration methods emphasize switching sequences, feeder reconfiguration, and operator decision support. More recent work incorporates distributed energy resources, microgrid islanding, cold-load pickup, and multi-period restoration constraints. For example, two-stage restoration formulations separate the selection of a restored configuration from the sequencing of switching operations, allowing distribution management systems to account for feasibility and operational timing \cite{poudel2020twoStage}. Other restoration studies examine how DER placement and operating assumptions change restoration order and unserved-energy outcomes in test feeders such as IEEE 123 \cite{rhodes2021derRestoration}.

These optimization-based methods provide important physical and operational foundations, but they can be computationally challenging for fast online use when the feeder is large, unbalanced, and subject to many possible outage scenarios. Learning-based controllers address this online speed challenge by shifting computation into an offline training phase. However, most restoration optimization and learning studies focus on nominal performance, uncertainty, or resilience to physical outages. They generally do not audit how a cyber attacker can manipulate the observation supplied to the restoration decision maker while leaving the physical feeder unchanged. This paper keeps the physical restoration evaluation of prior work, but places the cyber observation channel at the center of the threat model.

\subsection{Learning-Based and Graph-Based Restoration Control}

Reinforcement learning has been applied to sequential power-system and microgrid restoration because restoration naturally fits a Markov decision process: the controller observes an outage state, selects an action, receives a reward based on restored load and constraint satisfaction, and repeats over a restoration horizon. Deep reinforcement learning has been used for resilient load restoration and resource coordination in microgrids with mobile or distributed energy resources \cite{yao2019drlRestoration}. These methods show the promise of learned policies for online restoration decisions, especially when the action space is complex or when a fast response is desired.

Graph-based learning is especially relevant for distribution feeders because electrical networks are sparse graphs whose topology changes with switching actions. Graph neural networks and graph-network models provide a way to process node features, edge features, and global graph features while sharing parameters across network components \cite{kipf2017gcn,battaglia2018graphNetworks}. In restoration control, this structure can represent bus voltages, load states, line or switch features, branch-flow proxies, topology adjacency, and global restoration indicators. A GRL controller can therefore exploit feeder topology more directly than a flat multilayer perceptron.

The security implication is that the learned policy's observation vector is heterogeneous. It contains continuous measurements, binary status or availability features, graph connectivity information, and global scalar features. A cyber attack against such a controller need not look like a classical analog measurement attack. Corrupting an action mask can remove an otherwise valuable load-restoration action from the policy's feasible set; corrupting topology status can change the graph over which messages are passed; corrupting branch-flow or voltage measurements can bias the policy's perception of feasibility or loading. This paper therefore evaluates the controller observation as a structured attack surface rather than treating the learned controller as a black-box scalar input-output map.

\subsection{False-Data Injection Attacks in Power Systems}

The canonical FDIA literature shows that an attacker with sufficient system knowledge can construct measurement corruptions that evade conventional residual-based bad-data detection in power-system state estimation \cite{liu2009false}. Subsequent work studied the operational consequences of malicious data, including state-estimation error, market manipulation, and physical line-flow effects \cite{kosut2011malicious,xie2011false}. Defense work has examined meter protection, graph-theoretic measurement hardening, and attack localization strategies \cite{bi2014graphical}.

This paper is closely related to that literature in spirit but differs in the controlled object and impact metric. We do not study whether a state estimator is biased. Instead, we ask whether a restoration controller chooses a different action when its observation is corrupted, and whether that action reduces restored energy when evaluated on the true feeder. This distinction changes which observations matter. In state estimation, voltage, current, power injection, and power-flow measurements are the primary attack channels. In learned restoration, action/status masks and topology/status observations can be equally or more important because they directly determine the candidate action set and the graph representation consumed by the policy.

The paper also differs from stealth-only FDIA studies. A stealthy attack is important only if it creates operational harm. We therefore report both detection outcomes and physical degradation. We separately track accidental improvements, damaging-case counts, residual degradation after mitigation, and recovery ratios. This accounting avoids overstating risk when an attack changes measurements but does not reduce restored energy.

\subsection{Adversarial Reinforcement Learning}

Adversarial machine learning and adversarial reinforcement learning have shown that neural policies can be vulnerable to small perturbations of their observations \cite{huang2017adversarialPolicies,zhang2020robustDRL}. Those studies are important because they establish that learned policies may be brittle at decision time. However, generic adversarial-RL benchmarks often optimize perturbations against reward in simulated tasks without a domain-specific cyber-physical measurement model. Power-system restoration requires additional structure: the observation channels have engineering meaning, the selected action must be evaluated through power-flow physics, and not every numerical perturbation is a plausible cyber attack.

Our work adopts the adversarial-RL viewpoint that test-time observation manipulation can change a learned policy's action, but constrains the attack to named controller observations in a distribution restoration environment. This gives each attack a physical or operational interpretation, such as a false-unavailable load action, a false-open topology edge, a false-low voltage measurement, or a hidden branch-flow observation. The result is an attack assessment that is both policy-aware and power-system interpretable.

\subsection{Detection and Mitigation for Learning-Enabled Restoration}

Bad-data detection, residual checks, and measurement hardening are standard tools in power-system cybersecurity, but learning-enabled restoration adds a second layer: even when an attack is detected, the system must still choose a restoration action. A detector can identify that an observation is suspicious, but a poor mitigation strategy may still select an action that leaves substantial load unrestored. Conversely, a robust candidate-screening strategy can recover performance even when the attacked policy action is poor.

For this reason, this paper separates detection from mitigation. Detection profiles include residual-based checks and status-integrity checks, with relaxed profiles used to evaluate bypass risk. Mitigation then screens candidate actions, including attacked-observation candidates, detector-mitigated candidates, trusted-reference candidates, and horizon-screened candidates. This separation lets the paper identify whether residual risk comes from missed detection or from weak action selection after detection. The IEEE 34 and IEEE 123 case studies show that both effects matter: relaxed status assumptions can leave high-impact action/status attacks unresolved in IEEE 34, while IEEE 123 hard cases are primarily mitigation-selection failures after detection.

\subsection{Position of This Paper}

In summary, prior distribution restoration work provides optimization and learning foundations, FDIA work provides cyber-measurement attack models, and adversarial-RL work shows that learned policies can be sensitive to observation perturbations. The gap addressed here is the intersection of these areas: an auditable, observation-level, cyber-physical FDIA assessment for GRL-based distribution restoration. The distinguishing features are (i) named controller-observation attack registries, (ii) true-physics OpenDSS evaluation of every clean, attacked, and mitigated action, (iii) restored-energy degradation as the impact metric, (iv) cross-feeder IEEE 34 and IEEE 123 evidence, and (v) explicit separation of detection performance from mitigation residual risk.


\section{Introduction}
\IEEEPARstart{D}{istribution} system restoration is a time-critical control problem in which operators and automation systems must reconnect loads, reconfigure feeders, and coordinate available switching or controllable load resources after faults and outages. Modern active distribution networks increase the value of automated restoration because they include remotely controllable switches, distributed energy resources, microgrid-capable regions, and increasingly detailed sensing. At the same time, these capabilities expand the cyber-physical attack surface: restoration decisions depend not only on the true electrical state of the feeder, but also on the cyber observation made available to the restoration controller.

Learning-based restoration controllers are attractive in this setting because they can produce fast decisions after training and can naturally represent sequential restoration objectives. Graph reinforcement learning (GRL), in particular, is well matched to distribution feeders because feeder topology, switch state, voltage, branch-flow, and load connectivity are graph-structured quantities. A GRL policy can therefore use node, edge, and global features to map post-outage observations to switching and load-restoration actions. However, the same observation-rich structure that makes GRL useful also creates a security concern: if an attacker corrupts selected entries of the controller observation, the policy may choose a physically poor restoration action even though the underlying feeder has not changed.

False-data injection attacks (FDIAs) have been studied extensively in power-system monitoring and state estimation, where carefully designed measurement corruptions can bypass conventional bad-data detection and bias estimated system states \cite{liu2009false,kosut2011malicious,xie2011false,bi2014graphical}. Those results show that cyber measurements cannot be treated as automatically trustworthy. Yet restoration control differs from state estimation in several important ways. The consequence of an attack is not only an erroneous estimate; it is a changed restoration action that is then applied to the physical feeder. Moreover, the most security-critical observation may not be a traditional voltage or flow measurement. It may be a binary action mask, a topology/status input, or another controller-specific feature that constrains the learned policy's action selection.

This paper studies observation-level FDIAs against learning-based distribution restoration under a strict cyber-physical evaluation rule. The attacker corrupts only the observation received by the restoration controller. The physical feeder, outage scenario, and OpenDSS evaluation model remain the ground truth. For each deterministic outage-edge scenario, we compare the clean controller action, the attacked controller action, and the mitigated action on the same physical feeder state. Attack impact is measured using positive restored-energy degradation relative to the clean baseline, so cases in which an attack accidentally improves the restoration outcome are not counted as damage.

The study focuses on direct observations available to the current restoration controller: action/status masks, topology/status inputs, voltage measurements, branch-flow measurements, and global restoration scalars. This boundary is important. We do not claim direct attacks on load demand, DER availability, or explicit fault-status observations unless those quantities are exposed in the controller observation model. Instead, the goal is to quantify which existing controller-observation channels can create physical degradation when corrupted.

The paper is built on two frozen case-study evidence packages. The IEEE 34-bus study shows that action/status attacks dominate the observed vulnerability. In particular, false-unavailable dispatch-load attacks centered on the `Load\_s890' action/status observation create multi-MWh restored-energy degradation, and coordinated attacks involving the same observation reach mean physical degradation above 3 MWh. The IEEE 123-bus study shows a different vulnerability structure. In the baseline family study, branch-flow hiding is the strongest aggregate attack family, while exhaustive action-mask criticality reveals individual load-control observations whose impact is much larger than the baseline action-mask family average. These cross-feeder differences demonstrate that restoration security cannot be summarized by one generic FDIA family; it requires registry-driven criticality analysis over the actual observations used by the controller.

We also evaluate detection and mitigation. The detection layer includes residual and status-integrity checks, with strict and relaxed profiles used to quantify robustness and bypass risk. The mitigation layer screens candidate restoration actions after detection. In IEEE 34, trusted-reference screening substantially reduces residual degradation for high-impact action/status attacks, while relaxed status assumptions with weak candidate screening can leave large residual losses. In IEEE 123, the completed defense-stress study finds no high-impact detection bypass cases, but identifies detected low-recovery residual cases; advanced mitigation variants based on detector-selected, trusted-reference, and horizon-screened candidate actions remove the residual degradation in those hard cases.

The main contributions of this paper are:
\begin{itemize}
    \item An auditable observation-level FDIA assessment workflow for learning-based distribution restoration, where cyber observations are corrupted but all clean, attacked, and mitigated actions are evaluated on the true OpenDSS feeder state.
    \item A controller-observation attack taxonomy covering action/status masks, topology/status inputs, voltage measurements, branch-flow measurements, and global restoration scalars, with explicit claim boundaries for observations not exposed to the current controller.
    \item Frozen IEEE 34-bus and IEEE 123-bus GRL case studies with deterministic outage-edge replay, restored-energy degradation metrics, damaging-case counts, detection rates, bypass rates, mitigation recovery, and residual-risk accounting.
    \item Cross-feeder evidence that the dominant vulnerability mode is feeder and controller-context dependent: IEEE 34 is driven by dispatch-load action/status observations, whereas IEEE 123 combines branch-flow family effects with high-impact individual action-mask targets.
    \item A defense evaluation showing that strict status-integrity detection and trusted-reference or horizon-screened candidate mitigation can substantially reduce residual restoration loss, while detection alone is not sufficient when mitigation candidate selection remains weak.
\end{itemize}

The remainder of the paper is organized as follows. Section II reviews related work on distribution restoration, learning-based restoration, FDIA security, and adversarial reinforcement learning. Section III introduces the restoration environment, observation model, threat model, and defense model. Section IV describes the deterministic experimental workflow and evidence packages. Sections V and VI present the IEEE 34-bus and IEEE 123-bus results, respectively. Section VII discusses companion controller evidence and the boundary of the current claims. Section VIII concludes and outlines the planned IEEE 8500-bus HPC-scale extension.

\section{Related Work}

\subsection{Distribution System Restoration}

Distribution system restoration has traditionally been formulated as an optimization, search, or expert-system problem in which the objective is to maximize restored load or minimize unserved energy subject to radiality, voltage, thermal, switching, and operational constraints. Classical feeder restoration methods emphasize switching sequences, feeder reconfiguration, and operator decision support. More recent work incorporates distributed energy resources, microgrid islanding, cold-load pickup, and multi-period restoration constraints. For example, two-stage restoration formulations separate the selection of a restored configuration from the sequencing of switching operations, allowing distribution management systems to account for feasibility and operational timing \cite{poudel2020twoStage}. Other restoration studies examine how DER placement and operating assumptions change restoration order and unserved-energy outcomes in test feeders such as IEEE 123 \cite{rhodes2021derRestoration}.

These optimization-based methods provide important physical and operational foundations, but they can be computationally challenging for fast online use when the feeder is large, unbalanced, and subject to many possible outage scenarios. Learning-based controllers address this online speed challenge by shifting computation into an offline training phase. However, most restoration optimization and learning studies focus on nominal performance, uncertainty, or resilience to physical outages. They generally do not audit how a cyber attacker can manipulate the observation supplied to the restoration decision maker while leaving the physical feeder unchanged. This paper keeps the physical restoration evaluation of prior work, but places the cyber observation channel at the center of the threat model.

\subsection{Learning-Based and Graph-Based Restoration Control}

Reinforcement learning has been applied to sequential power-system and microgrid restoration because restoration naturally fits a Markov decision process: the controller observes an outage state, selects an action, receives a reward based on restored load and constraint satisfaction, and repeats over a restoration horizon. Deep reinforcement learning has been used for resilient load restoration and resource coordination in microgrids with mobile or distributed energy resources \cite{yao2019drlRestoration}. These methods show the promise of learned policies for online restoration decisions, especially when the action space is complex or when a fast response is desired.

Graph-based learning is especially relevant for distribution feeders because electrical networks are sparse graphs whose topology changes with switching actions. Graph neural networks and graph-network models provide a way to process node features, edge features, and global graph features while sharing parameters across network components \cite{kipf2017gcn,battaglia2018graphNetworks}. In restoration control, this structure can represent bus voltages, load states, line or switch features, branch-flow proxies, topology adjacency, and global restoration indicators. A GRL controller can therefore exploit feeder topology more directly than a flat multilayer perceptron.

The security implication is that the learned policy's observation vector is heterogeneous. It contains continuous measurements, binary status or availability features, graph connectivity information, and global scalar features. A cyber attack against such a controller need not look like a classical analog measurement attack. Corrupting an action mask can remove an otherwise valuable load-restoration action from the policy's feasible set; corrupting topology status can change the graph over which messages are passed; corrupting branch-flow or voltage measurements can bias the policy's perception of feasibility or loading. This paper therefore evaluates the controller observation as a structured attack surface rather than treating the learned controller as a black-box scalar input-output map.

\subsection{False-Data Injection Attacks in Power Systems}

The canonical FDIA literature shows that an attacker with sufficient system knowledge can construct measurement corruptions that evade conventional residual-based bad-data detection in power-system state estimation \cite{liu2009false}. Subsequent work studied the operational consequences of malicious data, including state-estimation error, market manipulation, and physical line-flow effects \cite{kosut2011malicious,xie2011false}. Defense work has examined meter protection, graph-theoretic measurement hardening, and attack localization strategies \cite{bi2014graphical}.

This paper is closely related to that literature in spirit but differs in the controlled object and impact metric. We do not study whether a state estimator is biased. Instead, we ask whether a restoration controller chooses a different action when its observation is corrupted, and whether that action reduces restored energy when evaluated on the true feeder. This distinction changes which observations matter. In state estimation, voltage, current, power injection, and power-flow measurements are the primary attack channels. In learned restoration, action/status masks and topology/status observations can be equally or more important because they directly determine the candidate action set and the graph representation consumed by the policy.

The paper also differs from stealth-only FDIA studies. A stealthy attack is important only if it creates operational harm. We therefore report both detection outcomes and physical degradation. We separately track accidental improvements, damaging-case counts, residual degradation after mitigation, and recovery ratios. This accounting avoids overstating risk when an attack changes measurements but does not reduce restored energy.

\subsection{Adversarial Reinforcement Learning}

Adversarial machine learning and adversarial reinforcement learning have shown that neural policies can be vulnerable to small perturbations of their observations \cite{huang2017adversarialPolicies,zhang2020robustDRL}. Those studies are important because they establish that learned policies may be brittle at decision time. However, generic adversarial-RL benchmarks often optimize perturbations against reward in simulated tasks without a domain-specific cyber-physical measurement model. Power-system restoration requires additional structure: the observation channels have engineering meaning, the selected action must be evaluated through power-flow physics, and not every numerical perturbation is a plausible cyber attack.

Our work adopts the adversarial-RL viewpoint that test-time observation manipulation can change a learned policy's action, but constrains the attack to named controller observations in a distribution restoration environment. This gives each attack a physical or operational interpretation, such as a false-unavailable load action, a false-open topology edge, a false-low voltage measurement, or a hidden branch-flow observation. The result is an attack assessment that is both policy-aware and power-system interpretable.

\subsection{Detection and Mitigation for Learning-Enabled Restoration}

Bad-data detection, residual checks, and measurement hardening are standard tools in power-system cybersecurity, but learning-enabled restoration adds a second layer: even when an attack is detected, the system must still choose a restoration action. A detector can identify that an observation is suspicious, but a poor mitigation strategy may still select an action that leaves substantial load unrestored. Conversely, a robust candidate-screening strategy can recover performance even when the attacked policy action is poor.

For this reason, this paper separates detection from mitigation. Detection profiles include residual-based checks and status-integrity checks, with relaxed profiles used to evaluate bypass risk. Mitigation then screens candidate actions, including attacked-observation candidates, detector-mitigated candidates, trusted-reference candidates, and horizon-screened candidates. This separation lets the paper identify whether residual risk comes from missed detection or from weak action selection after detection. The IEEE 34 and IEEE 123 case studies show that both effects matter: relaxed status assumptions can leave high-impact action/status attacks unresolved in IEEE 34, while IEEE 123 hard cases are primarily mitigation-selection failures after detection.

\subsection{Position of This Paper}

In summary, prior distribution restoration work provides optimization and learning foundations, FDIA work provides cyber-measurement attack models, and adversarial-RL work shows that learned policies can be sensitive to observation perturbations. The gap addressed here is the intersection of these areas: an auditable, observation-level, cyber-physical FDIA assessment for GRL-based distribution restoration. The distinguishing features are (i) named controller-observation attack registries, (ii) true-physics OpenDSS evaluation of every clean, attacked, and mitigated action, (iii) restored-energy degradation as the impact metric, (iv) cross-feeder IEEE 34 and IEEE 123 evidence, and (v) explicit separation of detection performance from mitigation residual risk.


\section{Introduction}
\IEEEPARstart{D}{istribution} system restoration is a time-critical control problem in which operators and automation systems must reconnect loads, reconfigure feeders, and coordinate available switching or controllable load resources after faults and outages. Modern active distribution networks increase the value of automated restoration because they include remotely controllable switches, distributed energy resources, microgrid-capable regions, and increasingly detailed sensing. At the same time, these capabilities expand the cyber-physical attack surface: restoration decisions depend not only on the true electrical state of the feeder, but also on the cyber observation made available to the restoration controller.

Learning-based restoration controllers are attractive in this setting because they can produce fast decisions after training and can naturally represent sequential restoration objectives. Graph reinforcement learning (GRL), in particular, is well matched to distribution feeders because feeder topology, switch state, voltage, branch-flow, and load connectivity are graph-structured quantities. A GRL policy can therefore use node, edge, and global features to map post-outage observations to switching and load-restoration actions. However, the same observation-rich structure that makes GRL useful also creates a security concern: if an attacker corrupts selected entries of the controller observation, the policy may choose a physically poor restoration action even though the underlying feeder has not changed.

False-data injection attacks (FDIAs) have been studied extensively in power-system monitoring and state estimation, where carefully designed measurement corruptions can bypass conventional bad-data detection and bias estimated system states \cite{liu2009false,kosut2011malicious,xie2011false,bi2014graphical}. Those results show that cyber measurements cannot be treated as automatically trustworthy. Yet restoration control differs from state estimation in several important ways. The consequence of an attack is not only an erroneous estimate; it is a changed restoration action that is then applied to the physical feeder. Moreover, the most security-critical observation may not be a traditional voltage or flow measurement. It may be a binary action mask, a topology/status input, or another controller-specific feature that constrains the learned policy's action selection.

This paper studies observation-level FDIAs against learning-based distribution restoration under a strict cyber-physical evaluation rule. The attacker corrupts only the observation received by the restoration controller. The physical feeder, outage scenario, and OpenDSS evaluation model remain the ground truth. For each deterministic outage-edge scenario, we compare the clean controller action, the attacked controller action, and the mitigated action on the same physical feeder state. Attack impact is measured using positive restored-energy degradation relative to the clean baseline, so cases in which an attack accidentally improves the restoration outcome are not counted as damage.

The study focuses on direct observations available to the current restoration controller: action/status masks, topology/status inputs, voltage measurements, branch-flow measurements, and global restoration scalars. This boundary is important. We do not claim direct attacks on load demand, DER availability, or explicit fault-status observations unless those quantities are exposed in the controller observation model. Instead, the goal is to quantify which existing controller-observation channels can create physical degradation when corrupted.

The paper is built on two frozen case-study evidence packages. The IEEE 34-bus study shows that action/status attacks dominate the observed vulnerability. In particular, false-unavailable dispatch-load attacks centered on the `Load\_s890' action/status observation create multi-MWh restored-energy degradation, and coordinated attacks involving the same observation reach mean physical degradation above 3 MWh. The IEEE 123-bus study shows a different vulnerability structure. In the baseline family study, branch-flow hiding is the strongest aggregate attack family, while exhaustive action-mask criticality reveals individual load-control observations whose impact is much larger than the baseline action-mask family average. These cross-feeder differences demonstrate that restoration security cannot be summarized by one generic FDIA family; it requires registry-driven criticality analysis over the actual observations used by the controller.

We also evaluate detection and mitigation. The detection layer includes residual and status-integrity checks, with strict and relaxed profiles used to quantify robustness and bypass risk. The mitigation layer screens candidate restoration actions after detection. In IEEE 34, trusted-reference screening substantially reduces residual degradation for high-impact action/status attacks, while relaxed status assumptions with weak candidate screening can leave large residual losses. In IEEE 123, the completed defense-stress study finds no high-impact detection bypass cases, but identifies detected low-recovery residual cases; advanced mitigation variants based on detector-selected, trusted-reference, and horizon-screened candidate actions remove the residual degradation in those hard cases.

The main contributions of this paper are:
\begin{itemize}
    \item An auditable observation-level FDIA assessment workflow for learning-based distribution restoration, where cyber observations are corrupted but all clean, attacked, and mitigated actions are evaluated on the true OpenDSS feeder state.
    \item A controller-observation attack taxonomy covering action/status masks, topology/status inputs, voltage measurements, branch-flow measurements, and global restoration scalars, with explicit claim boundaries for observations not exposed to the current controller.
    \item Frozen IEEE 34-bus and IEEE 123-bus GRL case studies with deterministic outage-edge replay, restored-energy degradation metrics, damaging-case counts, detection rates, bypass rates, mitigation recovery, and residual-risk accounting.
    \item Cross-feeder evidence that the dominant vulnerability mode is feeder and controller-context dependent: IEEE 34 is driven by dispatch-load action/status observations, whereas IEEE 123 combines branch-flow family effects with high-impact individual action-mask targets.
    \item A defense evaluation showing that strict status-integrity detection and trusted-reference or horizon-screened candidate mitigation can substantially reduce residual restoration loss, while detection alone is not sufficient when mitigation candidate selection remains weak.
\end{itemize}

The remainder of the paper is organized as follows. Section II reviews related work on distribution restoration, learning-based restoration, FDIA security, and adversarial reinforcement learning. Section III introduces the restoration environment, observation model, threat model, and defense model. Section IV describes the deterministic experimental workflow and evidence packages. Sections V and VI present the IEEE 34-bus and IEEE 123-bus results, respectively. Section VII discusses companion controller evidence and the boundary of the current claims. Section VIII concludes and outlines the planned IEEE 8500-bus HPC-scale extension.

\section{Related Work}

\subsection{Distribution System Restoration}

Distribution system restoration has traditionally been formulated as an optimization, search, or expert-system problem in which the objective is to maximize restored load or minimize unserved energy subject to radiality, voltage, thermal, switching, and operational constraints. Classical feeder restoration methods emphasize switching sequences, feeder reconfiguration, and operator decision support. More recent work incorporates distributed energy resources, microgrid islanding, cold-load pickup, and multi-period restoration constraints. For example, two-stage restoration formulations separate the selection of a restored configuration from the sequencing of switching operations, allowing distribution management systems to account for feasibility and operational timing \cite{poudel2020twoStage}. Other restoration studies examine how DER placement and operating assumptions change restoration order and unserved-energy outcomes in test feeders such as IEEE 123 \cite{rhodes2021derRestoration}.

These optimization-based methods provide important physical and operational foundations, but they can be computationally challenging for fast online use when the feeder is large, unbalanced, and subject to many possible outage scenarios. Learning-based controllers address this online speed challenge by shifting computation into an offline training phase. However, most restoration optimization and learning studies focus on nominal performance, uncertainty, or resilience to physical outages. They generally do not audit how a cyber attacker can manipulate the observation supplied to the restoration decision maker while leaving the physical feeder unchanged. This paper keeps the physical restoration evaluation of prior work, but places the cyber observation channel at the center of the threat model.

\subsection{Learning-Based and Graph-Based Restoration Control}

Reinforcement learning has been applied to sequential power-system and microgrid restoration because restoration naturally fits a Markov decision process: the controller observes an outage state, selects an action, receives a reward based on restored load and constraint satisfaction, and repeats over a restoration horizon. Deep reinforcement learning has been used for resilient load restoration and resource coordination in microgrids with mobile or distributed energy resources \cite{yao2019drlRestoration}. These methods show the promise of learned policies for online restoration decisions, especially when the action space is complex or when a fast response is desired.

Graph-based learning is especially relevant for distribution feeders because electrical networks are sparse graphs whose topology changes with switching actions. Graph neural networks and graph-network models provide a way to process node features, edge features, and global graph features while sharing parameters across network components \cite{kipf2017gcn,battaglia2018graphNetworks}. In restoration control, this structure can represent bus voltages, load states, line or switch features, branch-flow proxies, topology adjacency, and global restoration indicators. A GRL controller can therefore exploit feeder topology more directly than a flat multilayer perceptron.

The security implication is that the learned policy's observation vector is heterogeneous. It contains continuous measurements, binary status or availability features, graph connectivity information, and global scalar features. A cyber attack against such a controller need not look like a classical analog measurement attack. Corrupting an action mask can remove an otherwise valuable load-restoration action from the policy's feasible set; corrupting topology status can change the graph over which messages are passed; corrupting branch-flow or voltage measurements can bias the policy's perception of feasibility or loading. This paper therefore evaluates the controller observation as a structured attack surface rather than treating the learned controller as a black-box scalar input-output map.

\subsection{False-Data Injection Attacks in Power Systems}

The canonical FDIA literature shows that an attacker with sufficient system knowledge can construct measurement corruptions that evade conventional residual-based bad-data detection in power-system state estimation \cite{liu2009false}. Subsequent work studied the operational consequences of malicious data, including state-estimation error, market manipulation, and physical line-flow effects \cite{kosut2011malicious,xie2011false}. Defense work has examined meter protection, graph-theoretic measurement hardening, and attack localization strategies \cite{bi2014graphical}.

This paper is closely related to that literature in spirit but differs in the controlled object and impact metric. We do not study whether a state estimator is biased. Instead, we ask whether a restoration controller chooses a different action when its observation is corrupted, and whether that action reduces restored energy when evaluated on the true feeder. This distinction changes which observations matter. In state estimation, voltage, current, power injection, and power-flow measurements are the primary attack channels. In learned restoration, action/status masks and topology/status observations can be equally or more important because they directly determine the candidate action set and the graph representation consumed by the policy.

The paper also differs from stealth-only FDIA studies. A stealthy attack is important only if it creates operational harm. We therefore report both detection outcomes and physical degradation. We separately track accidental improvements, damaging-case counts, residual degradation after mitigation, and recovery ratios. This accounting avoids overstating risk when an attack changes measurements but does not reduce restored energy.

\subsection{Adversarial Reinforcement Learning}

Adversarial machine learning and adversarial reinforcement learning have shown that neural policies can be vulnerable to small perturbations of their observations \cite{huang2017adversarialPolicies,zhang2020robustDRL}. Those studies are important because they establish that learned policies may be brittle at decision time. However, generic adversarial-RL benchmarks often optimize perturbations against reward in simulated tasks without a domain-specific cyber-physical measurement model. Power-system restoration requires additional structure: the observation channels have engineering meaning, the selected action must be evaluated through power-flow physics, and not every numerical perturbation is a plausible cyber attack.

Our work adopts the adversarial-RL viewpoint that test-time observation manipulation can change a learned policy's action, but constrains the attack to named controller observations in a distribution restoration environment. This gives each attack a physical or operational interpretation, such as a false-unavailable load action, a false-open topology edge, a false-low voltage measurement, or a hidden branch-flow observation. The result is an attack assessment that is both policy-aware and power-system interpretable.

\subsection{Detection and Mitigation for Learning-Enabled Restoration}

Bad-data detection, residual checks, and measurement hardening are standard tools in power-system cybersecurity, but learning-enabled restoration adds a second layer: even when an attack is detected, the system must still choose a restoration action. A detector can identify that an observation is suspicious, but a poor mitigation strategy may still select an action that leaves substantial load unrestored. Conversely, a robust candidate-screening strategy can recover performance even when the attacked policy action is poor.

For this reason, this paper separates detection from mitigation. Detection profiles include residual-based checks and status-integrity checks, with relaxed profiles used to evaluate bypass risk. Mitigation then screens candidate actions, including attacked-observation candidates, detector-mitigated candidates, trusted-reference candidates, and horizon-screened candidates. This separation lets the paper identify whether residual risk comes from missed detection or from weak action selection after detection. The IEEE 34 and IEEE 123 case studies show that both effects matter: relaxed status assumptions can leave high-impact action/status attacks unresolved in IEEE 34, while IEEE 123 hard cases are primarily mitigation-selection failures after detection.

\subsection{Position of This Paper}

In summary, prior distribution restoration work provides optimization and learning foundations, FDIA work provides cyber-measurement attack models, and adversarial-RL work shows that learned policies can be sensitive to observation perturbations. The gap addressed here is the intersection of these areas: an auditable, observation-level, cyber-physical FDIA assessment for GRL-based distribution restoration. The distinguishing features are (i) named controller-observation attack registries, (ii) true-physics OpenDSS evaluation of every clean, attacked, and mitigated action, (iii) restored-energy degradation as the impact metric, (iv) cross-feeder IEEE 34 and IEEE 123 evidence, and (v) explicit separation of detection performance from mitigation residual risk.


\section{Introduction}
\IEEEPARstart{D}{istribution} system restoration is a time-critical control problem in which operators and automation systems must reconnect loads, reconfigure feeders, and coordinate available switching or controllable load resources after faults and outages. Modern active distribution networks increase the value of automated restoration because they include remotely controllable switches, distributed energy resources, microgrid-capable regions, and increasingly detailed sensing. At the same time, these capabilities expand the cyber-physical attack surface: restoration decisions depend not only on the true electrical state of the feeder, but also on the cyber observation made available to the restoration controller.

Learning-based restoration controllers are attractive in this setting because they can produce fast decisions after training and can naturally represent sequential restoration objectives. Graph reinforcement learning (GRL), in particular, is well matched to distribution feeders because feeder topology, switch state, voltage, branch-flow, and load connectivity are graph-structured quantities. A GRL policy can therefore use node, edge, and global features to map post-outage observations to switching and load-restoration actions. However, the same observation-rich structure that makes GRL useful also creates a security concern: if an attacker corrupts selected entries of the controller observation, the policy may choose a physically poor restoration action even though the underlying feeder has not changed.

False-data injection attacks (FDIAs) have been studied extensively in power-system monitoring and state estimation, where carefully designed measurement corruptions can bypass conventional bad-data detection and bias estimated system states \cite{liu2009false,kosut2011malicious,xie2011false,bi2014graphical}. Those results show that cyber measurements cannot be treated as automatically trustworthy. Yet restoration control differs from state estimation in several important ways. The consequence of an attack is not only an erroneous estimate; it is a changed restoration action that is then applied to the physical feeder. Moreover, the most security-critical observation may not be a traditional voltage or flow measurement. It may be a binary action mask, a topology/status input, or another controller-specific feature that constrains the learned policy's action selection.

This paper studies observation-level FDIAs against learning-based distribution restoration under a strict cyber-physical evaluation rule. The attacker corrupts only the observation received by the restoration controller. The physical feeder, outage scenario, and OpenDSS evaluation model remain the ground truth. For each deterministic outage-edge scenario, we compare the clean controller action, the attacked controller action, and the mitigated action on the same physical feeder state. Attack impact is measured using positive restored-energy degradation relative to the clean baseline, so cases in which an attack accidentally improves the restoration outcome are not counted as damage.

The study focuses on direct observations available to the current restoration controller: action/status masks, topology/status inputs, voltage measurements, branch-flow measurements, and global restoration scalars. This boundary is important. We do not claim direct attacks on load demand, DER availability, or explicit fault-status observations unless those quantities are exposed in the controller observation model. Instead, the goal is to quantify which existing controller-observation channels can create physical degradation when corrupted.

The paper is built on two frozen case-study evidence packages. The IEEE 34-bus study shows that action/status attacks dominate the observed vulnerability. In particular, false-unavailable dispatch-load attacks centered on the `Load\_s890' action/status observation create multi-MWh restored-energy degradation, and coordinated attacks involving the same observation reach mean physical degradation above 3 MWh. The IEEE 123-bus study shows a different vulnerability structure. In the baseline family study, branch-flow hiding is the strongest aggregate attack family, while exhaustive action-mask criticality reveals individual load-control observations whose impact is much larger than the baseline action-mask family average. These cross-feeder differences demonstrate that restoration security cannot be summarized by one generic FDIA family; it requires registry-driven criticality analysis over the actual observations used by the controller.

We also evaluate detection and mitigation. The detection layer includes residual and status-integrity checks, with strict and relaxed profiles used to quantify robustness and bypass risk. The mitigation layer screens candidate restoration actions after detection. In IEEE 34, trusted-reference screening substantially reduces residual degradation for high-impact action/status attacks, while relaxed status assumptions with weak candidate screening can leave large residual losses. In IEEE 123, the completed defense-stress study finds no high-impact detection bypass cases, but identifies detected low-recovery residual cases; advanced mitigation variants based on detector-selected, trusted-reference, and horizon-screened candidate actions remove the residual degradation in those hard cases.

The main contributions of this paper are:
\begin{itemize}
    \item An auditable observation-level FDIA assessment workflow for learning-based distribution restoration, where cyber observations are corrupted but all clean, attacked, and mitigated actions are evaluated on the true OpenDSS feeder state.
    \item A controller-observation attack taxonomy covering action/status masks, topology/status inputs, voltage measurements, branch-flow measurements, and global restoration scalars, with explicit claim boundaries for observations not exposed to the current controller.
    \item Frozen IEEE 34-bus and IEEE 123-bus GRL case studies with deterministic outage-edge replay, restored-energy degradation metrics, damaging-case counts, detection rates, bypass rates, mitigation recovery, and residual-risk accounting.
    \item Cross-feeder evidence that the dominant vulnerability mode is feeder and controller-context dependent: IEEE 34 is driven by dispatch-load action/status observations, whereas IEEE 123 combines branch-flow family effects with high-impact individual action-mask targets.
    \item A defense evaluation showing that strict status-integrity detection and trusted-reference or horizon-screened candidate mitigation can substantially reduce residual restoration loss, while detection alone is not sufficient when mitigation candidate selection remains weak.
\end{itemize}

The remainder of the paper is organized as follows. Section II reviews related work on distribution restoration, learning-based restoration, FDIA security, and adversarial reinforcement learning. Section III introduces the restoration environment, observation model, threat model, and defense model. Section IV describes the deterministic experimental workflow and evidence packages. Sections V and VI present the IEEE 34-bus and IEEE 123-bus results, respectively. Section VII discusses companion controller evidence and the boundary of the current claims. Section VIII concludes and outlines the planned IEEE 8500-bus HPC-scale extension.

\section{Related Work}

\subsection{Distribution System Restoration}

Distribution system restoration has traditionally been formulated as an optimization, search, or expert-system problem in which the objective is to maximize restored load or minimize unserved energy subject to radiality, voltage, thermal, switching, and operational constraints. Classical feeder restoration methods emphasize switching sequences, feeder reconfiguration, and operator decision support. More recent work incorporates distributed energy resources, microgrid islanding, cold-load pickup, and multi-period restoration constraints. For example, two-stage restoration formulations separate the selection of a restored configuration from the sequencing of switching operations, allowing distribution management systems to account for feasibility and operational timing \cite{poudel2020twoStage}. Other restoration studies examine how DER placement and operating assumptions change restoration order and unserved-energy outcomes in test feeders such as IEEE 123 \cite{rhodes2021derRestoration}.

These optimization-based methods provide important physical and operational foundations, but they can be computationally challenging for fast online use when the feeder is large, unbalanced, and subject to many possible outage scenarios. Learning-based controllers address this online speed challenge by shifting computation into an offline training phase. However, most restoration optimization and learning studies focus on nominal performance, uncertainty, or resilience to physical outages. They generally do not audit how a cyber attacker can manipulate the observation supplied to the restoration decision maker while leaving the physical feeder unchanged. This paper keeps the physical restoration evaluation of prior work, but places the cyber observation channel at the center of the threat model.

\subsection{Learning-Based and Graph-Based Restoration Control}

Reinforcement learning has been applied to sequential power-system and microgrid restoration because restoration naturally fits a Markov decision process: the controller observes an outage state, selects an action, receives a reward based on restored load and constraint satisfaction, and repeats over a restoration horizon. Deep reinforcement learning has been used for resilient load restoration and resource coordination in microgrids with mobile or distributed energy resources \cite{yao2019drlRestoration}. These methods show the promise of learned policies for online restoration decisions, especially when the action space is complex or when a fast response is desired.

Graph-based learning is especially relevant for distribution feeders because electrical networks are sparse graphs whose topology changes with switching actions. Graph neural networks and graph-network models provide a way to process node features, edge features, and global graph features while sharing parameters across network components \cite{kipf2017gcn,battaglia2018graphNetworks}. In restoration control, this structure can represent bus voltages, load states, line or switch features, branch-flow proxies, topology adjacency, and global restoration indicators. A GRL controller can therefore exploit feeder topology more directly than a flat multilayer perceptron.

The security implication is that the learned policy's observation vector is heterogeneous. It contains continuous measurements, binary status or availability features, graph connectivity information, and global scalar features. A cyber attack against such a controller need not look like a classical analog measurement attack. Corrupting an action mask can remove an otherwise valuable load-restoration action from the policy's feasible set; corrupting topology status can change the graph over which messages are passed; corrupting branch-flow or voltage measurements can bias the policy's perception of feasibility or loading. This paper therefore evaluates the controller observation as a structured attack surface rather than treating the learned controller as a black-box scalar input-output map.

\subsection{False-Data Injection Attacks in Power Systems}

The canonical FDIA literature shows that an attacker with sufficient system knowledge can construct measurement corruptions that evade conventional residual-based bad-data detection in power-system state estimation \cite{liu2009false}. Subsequent work studied the operational consequences of malicious data, including state-estimation error, market manipulation, and physical line-flow effects \cite{kosut2011malicious,xie2011false}. Defense work has examined meter protection, graph-theoretic measurement hardening, and attack localization strategies \cite{bi2014graphical}.

This paper is closely related to that literature in spirit but differs in the controlled object and impact metric. We do not study whether a state estimator is biased. Instead, we ask whether a restoration controller chooses a different action when its observation is corrupted, and whether that action reduces restored energy when evaluated on the true feeder. This distinction changes which observations matter. In state estimation, voltage, current, power injection, and power-flow measurements are the primary attack channels. In learned restoration, action/status masks and topology/status observations can be equally or more important because they directly determine the candidate action set and the graph representation consumed by the policy.

The paper also differs from stealth-only FDIA studies. A stealthy attack is important only if it creates operational harm. We therefore report both detection outcomes and physical degradation. We separately track accidental improvements, damaging-case counts, residual degradation after mitigation, and recovery ratios. This accounting avoids overstating risk when an attack changes measurements but does not reduce restored energy.

\subsection{Adversarial Reinforcement Learning}

Adversarial machine learning and adversarial reinforcement learning have shown that neural policies can be vulnerable to small perturbations of their observations \cite{huang2017adversarialPolicies,zhang2020robustDRL}. Those studies are important because they establish that learned policies may be brittle at decision time. However, generic adversarial-RL benchmarks often optimize perturbations against reward in simulated tasks without a domain-specific cyber-physical measurement model. Power-system restoration requires additional structure: the observation channels have engineering meaning, the selected action must be evaluated through power-flow physics, and not every numerical perturbation is a plausible cyber attack.

Our work adopts the adversarial-RL viewpoint that test-time observation manipulation can change a learned policy's action, but constrains the attack to named controller observations in a distribution restoration environment. This gives each attack a physical or operational interpretation, such as a false-unavailable load action, a false-open topology edge, a false-low voltage measurement, or a hidden branch-flow observation. The result is an attack assessment that is both policy-aware and power-system interpretable.

\subsection{Detection and Mitigation for Learning-Enabled Restoration}

Bad-data detection, residual checks, and measurement hardening are standard tools in power-system cybersecurity, but learning-enabled restoration adds a second layer: even when an attack is detected, the system must still choose a restoration action. A detector can identify that an observation is suspicious, but a poor mitigation strategy may still select an action that leaves substantial load unrestored. Conversely, a robust candidate-screening strategy can recover performance even when the attacked policy action is poor.

For this reason, this paper separates detection from mitigation. Detection profiles include residual-based checks and status-integrity checks, with relaxed profiles used to evaluate bypass risk. Mitigation then screens candidate actions, including attacked-observation candidates, detector-mitigated candidates, trusted-reference candidates, and horizon-screened candidates. This separation lets the paper identify whether residual risk comes from missed detection or from weak action selection after detection. The IEEE 34 and IEEE 123 case studies show that both effects matter: relaxed status assumptions can leave high-impact action/status attacks unresolved in IEEE 34, while IEEE 123 hard cases are primarily mitigation-selection failures after detection.

\subsection{Position of This Paper}

In summary, prior distribution restoration work provides optimization and learning foundations, FDIA work provides cyber-measurement attack models, and adversarial-RL work shows that learned policies can be sensitive to observation perturbations. The gap addressed here is the intersection of these areas: an auditable, observation-level, cyber-physical FDIA assessment for GRL-based distribution restoration. The distinguishing features are (i) named controller-observation attack registries, (ii) true-physics OpenDSS evaluation of every clean, attacked, and mitigated action, (iii) restored-energy degradation as the impact metric, (iv) cross-feeder IEEE 34 and IEEE 123 evidence, and (v) explicit separation of detection performance from mitigation residual risk.
